% ============================================================================
\section{The Automation Gap: Colombia Near Zero}
\label{sec:automation_gap}
% ============================================================================

Colombia occupies an anomalous position in the global automation landscape. Despite possessing Latin America's third-largest economy and a manufacturing sector that contributes approximately 11\% of GDP, the country's adoption of industrial robotics and advanced automation technologies remains negligible by any international benchmark. This section documents the magnitude of the gap, disaggregates its physical and cognitive dimensions, and argues that the disparity creates a precarious equilibrium: one in which a sudden cost shock---whether regulatory or technological---could trigger rapid, disorderly substitution rather than the gradual adjustment observed in high-income economies.

\subsection{Industrial Robot Density}
\label{subsec:robot_density}

According to data from the International Federation of Robotics \cite{ifr2024}, Colombia had an estimated total stock of 149 industrial robots in 2017, implying a robot density of approximately 2--5 units per 10,000 manufacturing workers. This figure places Colombia not merely below the global average density of 162 units per 10,000 workers, but effectively at the floor of the distribution among countries with meaningful manufacturing output. For perspective, South Korea leads globally with a density of 1,012 robots per 10,000 manufacturing workers, followed by Germany (429), Japan (419), and the United States (295). Even within Latin America, Colombia trails Mexico ($\sim$40--50 per 10,000) and Brazil ($\sim$15--20 per 10,000) by substantial margins.

Figure~\ref{fig:robot_timeseries} displays the time-series evolution of robot installations across selected economies, illustrating the exponential growth trajectories in East Asia and the stagnation in Latin American economies. Figure~\ref{fig:robot_density} provides a cross-sectional comparison of robot densities, making the scale of Colombia's gap visually apparent: the country's density is roughly 0.5\% of Korea's and 1.2\% of Germany's.

\begin{figure}[H]
\centering
\includegraphics[width=0.85\textwidth]{fig2_robot_timeseries.png}
\caption{Annual industrial robot installations by country, 2010--2023. Source: IFR World Robotics 2024.}
\label{fig:robot_timeseries}
\end{figure}

\begin{figure}[H]
\centering
\includegraphics[width=0.85\textwidth]{fig5_robot_density_comparison.png}
\caption{Robot density per 10,000 manufacturing workers, latest available year. Source: IFR World Robotics 2024.}
\label{fig:robot_density}
\end{figure}

The near-zero robot density has two important analytical implications. First, it means that the displacement effects documented by \cite{acemoglu2020robots} for the United States---where each additional robot per 1,000 workers reduced the employment-to-population ratio by 0.2 percentage points---have yet to materialize in Colombian manufacturing through the physical automation channel. Second, it implies that Colombia's manufacturing competitiveness relies almost entirely on labor cost arbitrage rather than capital deepening, a strategy that becomes increasingly fragile as regional competitors such as Mexico accelerate their automation investments in response to nearshoring demand.

\subsection{R\&D Investment}
\label{subsec:rd_investment}

The robot density gap is undergirded by a chronic underinvestment in research and development. Colombia allocates approximately 0.30\% of GDP to R\&D, a figure that has remained essentially unchanged for over a decade. This level of investment is an order of magnitude below leading innovators: South Korea devotes 5.0\% of GDP to R\&D, the United States 3.5\%, Japan 3.4\%, and Germany 3.1\%. Even Brazil, Colombia's closest regional comparator in economic structure, invests 1.15\% of GDP---nearly four times the Colombian level.

The consequences of this underinvestment are visible at the firm level. The most recent wave of the \textit{Encuesta de Desarrollo e Innovaci\'on Tecnol\'ogica} (EDIT IX) reveals that 75.4\% of Colombian manufacturing firms reported undertaking \textit{no innovation activity whatsoever} during the survey period. This figure captures not merely the absence of frontier R\&D but the absence of even incremental process improvements, technology licensing, or organizational innovation. The result is a manufacturing sector in which the modal firm operates with unchanged production technology year after year, relying on low wages rather than productivity gains to maintain viability.

\subsection{Physical vs.\ Cognitive Automation}
\label{subsec:physical_cognitive}

The portrait of Colombia as an automation laggard requires qualification along one critical dimension: the distinction between physical automation (industrial robots, automated production lines, CNC machinery) and cognitive automation (software-based AI, machine learning, natural language processing, robotic process automation). While Colombia's physical automation gap is extreme and shows no signs of closing, cognitive automation---particularly in services---is advancing rapidly and, in certain sectors, at rates comparable to high-income economies.

The financial sector exemplifies this divergence. According to \cite{asobancaria2024}, 73\% of Colombian financial entities have adopted some form of artificial intelligence, predominantly in credit scoring, fraud detection, customer service chatbots, and regulatory compliance. Bancolombia, the country's largest bank, employs approximately 4,000 technology workers and has emerged as a regional leader in the adoption of AI-assisted development tools such as GitHub Copilot. Similarly, the business process outsourcing (BPO) sector---which employs over 600,000 workers---reports AI adoption rates of approximately 80\%, driven by the imperative to maintain cost competitiveness against Indian and Philippine competitors.

These sectoral adoption patterns are partially captured by the Latin American Artificial Intelligence Index (ILIA) 2024 \cite{ilia2024}, which ranks Colombia fifth in the region with a composite score of 52.64 out of 100, behind Chile (73.07), Brazil (69.30), Uruguay (64.98), and Argentina. The gap relative to Chile is particularly notable: it reflects not differences in private-sector dynamism but rather the superior quality of Chile's AI governance frameworks, public data infrastructure, and university-industry linkages.

The coexistence of near-zero physical automation with rapidly advancing cognitive automation creates a distinctive risk profile. In the manufacturing sector, the binding constraint on automation is capital investment, and low robot density implies that displacement pressures will emerge gradually. In services, however, the binding constraint is software deployment---which requires minimal capital expenditure and can scale instantaneously across an organization. This asymmetry suggests that the most acute displacement risks in Colombia over the next decade will concentrate not in manufacturing but in formal services: precisely the sectors where the labor reform's cost increases would bite hardest.


% ============================================================================
\section{International Evidence: Labor Costs as Catalyst}
\label{sec:international}
% ============================================================================

The relationship between labor costs and automation adoption is not merely a theoretical proposition; it has been tested, in varied institutional contexts, across the world's leading manufacturing economies. This section reviews the evidence from six national experiences, selected to span the relevant dimensions of variation: wage levels, regulatory frameworks, demographic pressures, and industrial policy ambitions. In each case, we trace the mechanism by which labor cost dynamics---whether driven by minimum wage increases, social insurance mandates, or demographic wage pressure---have shaped the pace and character of automation adoption, with particular attention to the employment consequences and the policy responses that mediated them.

\subsection{South Korea}
\label{subsec:korea}

South Korea provides the starkest illustration of the labor cost--automation nexus. The country's robot density of 1,012 per 10,000 manufacturing workers---the highest in the world---did not emerge from a vacuum but was propelled by decades of rapid wage growth in a tightening labor market. The mechanism was tested acutely in 2018, when the Moon Jae-in government implemented a 16.4\% increase in the minimum wage in a single year. The result was a loss of approximately 170,000 manufacturing jobs, concentrated among small and medium enterprises that lacked the capital to absorb the cost increase and instead chose either automation or closure. Research by the Bank of Korea has estimated that an increase of one robot per 1,000 workers reduces the employment rate by 0.1 percentage points, an elasticity that, applied to Korea's density levels, implies substantial cumulative displacement.

The Korean experience is instructive for Colombia in two respects. First, it demonstrates that large, discrete increases in labor costs can trigger non-linear automation responses: firms that had been deferring automation investments because the cost differential was insufficient find that a regulatory shock pushes them past the threshold simultaneously, producing a wave of adoption rather than a gradual transition. Second, Korea's ability to manage the social consequences of this displacement---through a robust social insurance system, active labor market policies, and rapid job creation in new technology sectors---represents institutional capacities that Colombia conspicuously lacks.

\subsection{Germany}
\label{subsec:germany}

Germany offers a contrasting model: high robot density (429 per 10,000 manufacturing workers) coexisting with relatively stable manufacturing employment. Between 1993 and 2017, Germany lost only 19\% of its manufacturing employment share, compared to 33\% in the United States over the same period. Several institutional features explain this divergence. The Hartz labor market reforms of 2003--2005 reduced effective labor costs through the liberalization of temporary and part-time employment, the creation of ``mini-jobs,'' and the tightening of unemployment benefit conditionality. By reducing the labor cost wedge, the Hartz reforms paradoxically \textit{slowed} the automation incentive: with cheaper flexible labor available, the cost calculus favoring robot adoption became less compelling at the margin.

Germany's \textit{Industrie 4.0} strategy, launched in 2013, further illustrates the role of deliberate policy design. Rather than allowing automation to proceed as an uncoordinated firm-level response to cost pressures, the German government framed digitalization as a national competitiveness project, embedding it within the dual vocational training system (\textit{duale Ausbildung}) and the codetermination framework (\textit{Mitbestimmung}) that gives workers' councils a voice in technology adoption decisions. The result has been a pattern of automation that complements rather than substitutes for skilled labor, with productivity gains shared through collective bargaining. Colombia possesses none of these institutional buffers: no dual training system, no effective social dialogue on technology adoption, and no mechanism for workers to participate in firm-level automation decisions.

\subsection{China}
\label{subsec:china}

China's automation trajectory reflects the intersection of industrial policy ambition and demographic necessity. The \textit{Made in China 2025} initiative, announced in 2015, set explicit targets for robot density and domestic robot production, positioning automation as central to the country's transition from labor-intensive to technology-intensive manufacturing. Guangdong province alone committed \$135 billion to ``human-machine substitution'' programs that subsidized robot purchases for small and medium manufacturers. The results have been dramatic: China's robot density rose to 470 per 10,000 manufacturing workers, making it the third-most automated manufacturing economy globally.

The Chinese case demonstrates that automation can proceed rapidly when state subsidies reduce the effective capital cost of robots, even in an economy with relatively low wages. This finding is directly relevant to the \cite{acemoglu2020taxes} argument that tax and subsidy structures can distort automation decisions: China's subsidies artificially reduced the cost of capital relative to labor, accelerating substitution beyond the level that market conditions alone would have warranted. For Colombia, the Chinese experience raises an uncomfortable question: while Colombia's high parafiscal costs raise the price of labor, the country offers no countervailing subsidies to channel automation toward productivity-enhancing applications rather than pure labor substitution.

\subsection{Japan}
\label{subsec:japan}

Japan presents a demographically unique case. With 29\% of its population aged 65 or older and a working-age population that has been shrinking for two decades, Japan automates not primarily to \textit{eliminate} jobs but to \textit{fill} vacancies that the labor market can no longer supply. Robot density in Japan (419 per 10,000) reflects an economy in which automation is a response to labor scarcity rather than labor cost, and where the social acceptability of robots in the workplace is reinforced by cultural factors and the absence of immigration as a workforce supplement.

The Japanese model suggests that the employment consequences of automation are contingent on labor market tightness. In an economy with chronic labor surpluses---as in much of Colombia's informal sector---automation displaces workers into unemployment or deeper informality. In an economy with chronic labor shortages, automation absorbs tasks that would otherwise go unperformed, and the net employment effect may be positive through the reinstatement channel emphasized by \cite{acemoglu2019automation}. Colombia's labor market, with unemployment rates of 10--12\% and informality rates exceeding 55\%, bears no resemblance to Japan's, implying that automation in Colombia is far more likely to follow the displacement pathway.

\subsection{Mexico and Chile}
\label{subsec:mexico_chile}

Within Latin America, Mexico and Chile provide the most instructive comparisons. Mexico's robot density of approximately 40--50 per 10,000 manufacturing workers is modest by global standards but an order of magnitude above Colombia's. Critically, roughly 70\% of Mexico's installed robots are concentrated in the automotive sector, reflecting the integration of Mexican manufacturing into North American supply chains where automation standards are set by U.S. and German parent firms. The nearshoring wave triggered by U.S.--China trade tensions is now driving a new phase of automation investment in Mexico, as firms relocating production from Asia install state-of-the-art automated facilities rather than replicating the labor-intensive models of previous decades.

Chile offers a different lens: the automation of extractive industries. BHP's Escondida copper mine operates 33 autonomous haul trucks, while the BHP Spence mine has achieved fully autonomous operation. Critically, BHP reports that 80\% of the workers displaced by these autonomous systems were retrained and redeployed within the company, a figure enabled by Chile's relatively advanced vocational training infrastructure and BHP's own corporate commitments. The Chilean mining experience suggests that automation need not produce mass displacement when accompanied by robust reskilling programs---but also that such programs require institutional capacities and corporate cultures that are rare in Colombia's more fragmented industrial landscape.


% ============================================================================
\section{Results}
\label{sec:results}
% ============================================================================

This section presents the empirical results across four analytical levels: international panel regressions linking labor costs to robot adoption; individual-level models of automation risk in the Colombian labor market; firm-level evidence on the relationship between labor costs, capital intensity, and innovation; and a sectoral vulnerability analysis that integrates these dimensions into a composite index. Together, these results establish that Colombia faces a distinctive automation risk profile: low current adoption masking high latent vulnerability, concentrated among formal workers in cost-sensitive sectors, and poised to accelerate under plausible regulatory and technological scenarios.

\subsection{International Panel: Labor Costs and Robot Adoption}
\label{subsec:panel_results}

We estimate the determinants of robot density using an unbalanced panel of $N = 226$ country-year observations spanning 25 countries over the period 2010--2023. The panel includes all major robot-adopting economies as well as key emerging markets, with robot density data drawn from the IFR \cite{ifr2024} and macroeconomic controls from the Penn World Table 11.0 and World Development Indicators.

Table~\ref{tab:panel_results} reports the results from three specifications. The pooled OLS model (Column 1) yields an $R^2$ of 0.27, with the R\&D-to-GDP ratio entering with a positive and significant coefficient of 0.92 ($p < 0.01$), consistent with the hypothesis that innovation capacity is a prerequisite for robot adoption. However, the pooled OLS specification is subject to omitted variable bias from unobserved country-level heterogeneity in industrial structure, institutional quality, and automation culture.

\begin{table}[H]
\centering
\caption{Determinants of Robot Density: International Panel Regressions, 2010--2023}
\label{tab:panel_results}
\begin{threeparttable}
\begin{tabular}{lccc}
\toprule
& \textbf{Pooled OLS} & \textbf{Fixed Effects} & \textbf{Dynamic OLS} \\
& (1) & (2) & (3) \\
\midrule
$\log(\text{GDP per worker})$ & & $3.60^{***}$ & \\
& & $(0.82)$ & \\[4pt]
R\&D / GDP & $0.92^{***}$ & $-0.20^{***}$ & \\
& $(0.18)$ & $(0.06)$ & \\[4pt]
$\text{Robot density}_{t-1}$ & & & $1.01^{***}$ \\
& & & $(0.03)$ \\[4pt]
\midrule
$R^2$ & 0.27 & 0.62 & 0.97 \\
$N$ & 226 & 226 & 201 \\
Country FE & No & Yes & Yes \\
Year FE & No & Yes & Yes \\
\bottomrule
\end{tabular}
\begin{tablenotes}
\small
\item \textit{Notes:} Standard errors in parentheses. $^{***}p<0.01$, $^{**}p<0.05$, $^{*}p<0.10$. Dependent variable: robot density per 10,000 manufacturing workers. The Hausman test rejects the null of no systematic difference between FE and RE at the 1\% level, favoring the fixed effects specification.
\end{tablenotes}
\end{threeparttable}
\end{table}

The country fixed effects specification (Column 2) absorbs time-invariant heterogeneity and raises the $R^2$ to 0.62. In this model, the coefficient on log GDP per worker is 3.60 ($p < 0.001$), indicating that a 1\% increase in labor productivity is associated with a 3.6-unit increase in robot density---consistent with the complementarity between high-productivity production systems and automation. The R\&D-to-GDP coefficient switches sign to $-0.20$ ($p < 0.01$), reflecting the well-known within-country pattern in which R\&D intensity plateaus or declines as countries approach the technological frontier, even as their robot density continues to rise through deployment of existing technologies. A Hausman test decisively rejects the random effects specification in favor of fixed effects ($p < 0.001$).

The dynamic OLS specification (Column 3) introduces the lagged dependent variable, yielding a coefficient of 1.01 ($p < 0.001$) and an $R^2$ of 0.97. The near-unit autoregressive coefficient reveals extreme persistence in robot density: countries that have invested heavily in automation continue to do so, while laggards remain trapped in a low-automation equilibrium. This persistence is consistent with the theoretical prediction that automation generates self-reinforcing dynamics through learning-by-doing, network effects in the robotics ecosystem, and the accumulation of complementary human capital.

\begin{figure}[H]
\centering
\includegraphics[width=0.85\textwidth]{fig1_productivity_vs_robots.png}
\caption{Labor productivity (log GDP per worker) vs.\ robot density across 25 countries, 2010--2023. The fitted line reflects the fixed effects specification. Colombia is circled in the lower-left quadrant.}
\label{fig:productivity_robots}
\end{figure}

\begin{figure}[H]
\centering
\includegraphics[width=0.75\textwidth]{fig3_coefficient_plot.png}
\caption{Coefficient plot from the international panel fixed effects model. Bars represent 95\% confidence intervals.}
\label{fig:coefficient_plot}
\end{figure}

The most striking finding emerges from a counterfactual exercise: using the fixed effects model to predict Colombia's robot density based on its macroeconomic fundamentals---GDP per capita, labor productivity, manufacturing share, and R\&D intensity---the model predicts that Colombia should be installing approximately 1,200 or more robots per year. Actual annual installations range between 120 and 200 units, implying a gap of roughly one order of magnitude. This ``missing robots'' phenomenon cannot be explained by macroeconomic fundamentals alone; it reflects idiosyncratic barriers including the small average firm size in Colombian manufacturing, limited access to capital for technology investment, weak integration into global supply chains that mandate automation standards, and the availability of cheap informal labor as a substitute for capital investment.

\subsection{Automation Risk in the Colombian Labor Market}
\label{subsec:risk_colombia}

To characterize the distribution of automation risk across the Colombian workforce, we map the occupation-level automation probabilities estimated by \cite{frey2017future} to the occupational classification system used in the DANE's \textit{Gran Encuesta Integrada de Hogares} (GEIH) for 2024. The resulting dataset comprises $N = 111{,}672$ workers observed across four months, each assigned an automation probability based on their reported occupation at the one-digit level.

The aggregate distribution of automation risk is alarming. The mean automation probability across all workers is 0.489, with a median of 0.550, indicating that the distribution is slightly left-skewed but centered near the conventional threshold for ``high risk.'' Using the standard cutoff of $\Pr(\text{automation}) \geq 0.50$, we find that 53.8\% of the sampled workforce occupies high-risk occupations. This figure is broadly consistent with the estimates of \cite{morales2023risk}, who report 37\% of Colombian formal employment at high risk using a somewhat different mapping methodology, and with the 58\% upper bound estimated by \cite{fedesarrollo2023} for the Andean region.

To identify the individual and job characteristics associated with automation risk, we estimate a logistic regression of the form:

\begin{equation}
\label{eq:logit}
\Prob(\text{High Risk}_i = 1) = \Lambda\left(\beta_0 + \beta_1 \text{Formal}_i + \beta_2 \text{Female}_i + \beta_3 \log(\text{Income}_i) + \boldsymbol{\beta}_4' \mathbf{X}_i + \varepsilon_i\right)
\end{equation}

\noindent where $\Lambda(\cdot)$ is the logistic CDF and $\mathbf{X}_i$ includes education level (categorical, with \textit{secundaria} as the reference), age, firm size, and sector fixed effects. The model achieves a pseudo-$R^2$ of 0.308, an area under the ROC curve (AUC) of 0.851, and an out-of-sample classification accuracy of 79\%, indicating strong discriminatory power.

\begin{figure}[H]
\centering
\includegraphics[width=0.75\textwidth]{fig_roc_curve.png}
\caption{Receiver Operating Characteristic (ROC) curve for the logistic model of high automation risk. AUC $= 0.851$.}
\label{fig:roc_curve}
\end{figure}

Table~\ref{tab:odds_ratios} reports the key odds ratios. The results reveal a set of sharp asymmetries:

\begin{table}[H]
\centering
\caption{Logistic Regression: Odds Ratios for High Automation Risk ($\Pr \geq 0.50$)}
\label{tab:odds_ratios}
\begin{threeparttable}
\begin{tabular}{lcc}
\toprule
\textbf{Variable} & \textbf{Odds Ratio} & \textbf{95\% CI} \\
\midrule
Formal worker & $3.52^{***}$ & [3.21, 3.86] \\
Female & $0.67^{***}$ & [0.63, 0.71] \\
$\log(\text{Income})$ & $0.24^{***}$ & [0.22, 0.26] \\
Posgrado (vs.\ Secundaria) & $0.09^{***}$ & [0.07, 0.12] \\
Large firm ($>$200 employees) & $0.64^{***}$ & [0.58, 0.71] \\
\midrule
\multicolumn{3}{l}{\textit{Sector (vs.\ Comercio)}} \\
\quad Transport & $6.10^{***}$ & [5.31, 7.01] \\
\quad Mining & $7.86^{***}$ & [5.92, 10.43] \\
\quad Manufacturing & $3.51^{***}$ & [3.14, 3.92] \\
\bottomrule
\end{tabular}
\begin{tablenotes}
\small
\item \textit{Notes:} $N = 111{,}672$. Pseudo-$R^2 = 0.308$. AUC $= 0.851$. $^{***}p<0.01$. Reference categories: informal, male, secundaria, small firm ($<$10), comercio.
\end{tablenotes}
\end{threeparttable}
\end{table}

\begin{figure}[H]
\centering
\includegraphics[width=0.85\textwidth]{fig_forest_plot_odds_ratios.png}
\caption{Forest plot of odds ratios from the logistic model of high automation risk. Dots represent point estimates; horizontal lines represent 95\% confidence intervals.}
\label{fig:forest_plot}
\end{figure}

\textbf{Formality.} Formal workers face 3.52 times the odds of occupying a high-risk occupation compared to informal workers, controlling for all other covariates. This is the single most powerful predictor in the model and constitutes the core of the formality paradox discussed in Section~\ref{subsec:formality_paradox} below.

\textbf{Gender.} Female workers have 33\% lower odds of high automation risk (OR $= 0.67$), reflecting their concentration in service occupations---education, health care, personal services---that involve non-routine interpersonal tasks with lower automation potential. This finding does not imply that women are insulated from technological disruption; generative AI is increasingly capable of automating cognitive tasks in precisely the service sectors where women predominate, a point to which we return in Section~\ref{subsec:genai}.

\textbf{Income.} A one-unit increase in log income is associated with a 76\% reduction in the odds of high automation risk (OR $= 0.24$), consistent with the canonical finding that routine tasks---which are the most automatable---are concentrated in the middle and lower segments of the earnings distribution \cite{frey2017future}.

\textbf{Education.} The protective effect of education exhibits a steep gradient. Relative to workers with secondary education, those with postgraduate degrees face a 91\% reduction in the odds of high automation risk (OR $= 0.09$). This gradient is considerably steeper than those reported for high-income economies, reflecting the more polarized occupational structure of the Colombian labor market, where university-educated workers are disproportionately concentrated in managerial and professional roles with low automation potential.

\textbf{Firm size.} Workers in large firms (more than 200 employees) have 36\% lower odds of high automation risk (OR $= 0.64$). This result is initially counterintuitive---one might expect large firms to automate more aggressively---but reflects the fact that large Colombian firms are concentrated in sectors (financial services, professional services, telecommunications) with lower mean automation probabilities and more complex, non-routine task portfolios.

\textbf{Sectors.} The sectoral coefficients reveal extreme heterogeneity. Transport (OR $= 6.10$) and mining (OR $= 7.86$) exhibit the highest odds ratios, reflecting the dominance of routine physical tasks---driving, machine operation, extraction---that are highly susceptible to physical automation. Manufacturing (OR $= 3.51$) occupies an intermediate position.

\begin{figure}[H]
\centering
\includegraphics[width=0.85\textwidth]{fig_automation_risk_by_sector.png}
\caption{Mean automation probability by economic sector. The dashed line indicates the 0.50 high-risk threshold.}
\label{fig:risk_by_sector}
\end{figure}

Figure~\ref{fig:risk_by_sector} displays mean automation probabilities by sector. The highest-risk sectors are Transport (mean probability 0.656), Agriculture (0.614), and Households as employers (0.598). The lowest-risk sectors are Education (0.168), Professional services (0.234), and Health (0.286). The spread of nearly 50 percentage points between the most and least exposed sectors underscores the inadequacy of aggregate automation risk estimates: the Colombian labor market is not uniformly exposed but rather deeply segmented, with pockets of extreme vulnerability adjacent to sectors that face minimal near-term risk.

\begin{figure}[H]
\centering
\includegraphics[width=0.85\textwidth]{fig_automation_risk_by_education.png}
\caption{Distribution of automation risk by educational attainment. Each violin represents the conditional density of automation probabilities for the given education level.}
\label{fig:risk_by_education}
\end{figure}

\subsection{The Formality Paradox}
\label{subsec:formality_paradox}

The finding that formal workers face 3.52 times the odds of high automation risk demands careful interpretation, as it encodes two distinct mechanisms. The first is compositional: formal employment in Colombia is concentrated in manufacturing, transport, financial services, and large-scale retail---sectors with high shares of routine tasks. Informal employment, by contrast, is dominated by self-employment in personal services, street vending, construction day-labor, and small-scale agriculture---occupations that, while precarious, involve idiosyncratic, non-routine physical tasks that current technologies cannot easily replicate.

The second mechanism is economic: formal workers generate non-wage costs estimated at 46--58\% above base salary, including pension contributions (12\%), health insurance (8.5\%), occupational risk insurance (0.52--6.96\%), severance and interest on severance (9.3\%), vacation pay (4.2\%), and parafiscal contributions (4--9\%). Informal workers generate none of these costs. When a firm evaluates the cost-benefit calculus of automation, the relevant comparison is not the base wage versus the annualized cost of a robot, but the \textit{total} employer cost---including all mandatory contributions---versus the all-in cost of the technological substitute. The 52\% non-wage cost wedge dramatically shifts this calculus in favor of automation for every formal position.

\begin{figure}[H]
\centering
\includegraphics[width=0.85\textwidth]{fig_heatmap_sector_formality.png}
\caption{Heatmap of mean automation probability by sector and formality status. Darker shading indicates higher automation risk.}
\label{fig:heatmap_formality}
\end{figure}

Figure~\ref{fig:heatmap_formality} presents a heatmap of mean automation probabilities disaggregated by both sector and formality status. The pattern is nuanced. In manufacturing, informal workers face a higher \textit{technical} automation probability (0.66) than formal workers (0.55), reflecting the fact that informal manufacturing workers tend to perform more routine manual tasks. In agriculture, the gap is narrower: informal workers at 0.63 versus formal workers at 0.60. However, this comparison of technical probabilities is misleading if taken at face value. The \textit{effective} automation risk---the probability that a worker's job will actually be automated within a given time horizon---depends not only on the technical feasibility of automation but on the economic incentive to automate. And this incentive is vastly stronger for formal workers, whose total compensation cost is 46--58\% higher than their nominal wage.

We formalize this insight by estimating an interaction model in which the effect of formality on automation risk is allowed to vary by sector. The interaction terms are jointly significant ($p < 0.001$), confirming that the formality paradox is not a uniform phenomenon but varies in magnitude across the sectoral structure. The paradox is most acute in sectors where both technical automation potential and formality rates are high---precisely the sectors (manufacturing, transport, financial services) that account for the bulk of formal employment creation in Colombia.

The policy implication is stark: policies that increase the cost of formal employment---through higher parafiscal contributions, expanded benefits mandates, or reduced working hours at unchanged pay---do not merely reduce labor demand through the conventional channel of moving along a downward-sloping demand curve. They also shift firms toward the automation margin by widening the gap between total labor cost and the cost of technological substitutes. Informality, paradoxically, protects workers from automation precisely because it renders them invisible to the regulatory framework that generates the cost wedge incentivizing substitution.

\subsection{Firm-Level Evidence: Labor Costs, Capital Intensity, and Innovation}
\label{subsec:firm_results}

To move beyond occupation-level risk assessments and examine the actual behavioral responses of firms to labor cost pressures, we construct a matched firm-level dataset by merging the \textit{Encuesta Anual Manufacturera} (EAM) with the \textit{Encuesta de Desarrollo e Innovaci\'on Tecnol\'ogica} (EDIT). The merge, performed on establishment identifiers, yields 5,884 matched firms with a merge rate of 96.2\%, providing high confidence that the matched sample is representative of the manufacturing census.

We estimate a sequence of models relating labor cost measures to automation proxies:

\begin{table}[H]
\centering
\caption{Firm-Level Regressions: Labor Costs, Capital Intensity, and Innovation}
\label{tab:firm_results}
\begin{threeparttable}
\begin{tabular}{llccc}
\toprule
\textbf{Model} & \textbf{Specification} & $\boldsymbol{\hat{\beta}}$ & \textbf{Std.\ Err.} & $\boldsymbol{R^2}$ \\
\midrule
A & $\log(\text{Automation proxy})$ on $\log(\text{ULC})$ & $0.84^{***}$ & (0.04) & 0.43 \\
B & Investment rate on labor share of VA & $0.045^{***}$ & (0.008) & 0.11 \\
C & $\log(\text{Capital intensity})$ on $\log(\text{Labor cost/worker})$ & $1.27^{***}$ & (0.05) & 0.52 \\
\midrule
D & Logit: Innovator status on $\log(\text{ULC})$ & $<1^{**}$ & --- & --- \\
E2 & Investment intensity (conditional) on $\log(\text{ULC})$ & $0.81^{***}$ & (0.12) & 0.38 \\
\bottomrule
\end{tabular}
\begin{tablenotes}
\small
\item \textit{Notes:} $N = 5{,}884$ matched firms. $^{***}p<0.01$, $^{**}p<0.05$. Model A uses an automation proxy constructed from capital expenditure on machinery and equipment. Model D is a logistic regression where the dependent variable is a binary indicator for EDIT-classified innovator status.
\end{tablenotes}
\end{threeparttable}
\end{table}

\begin{figure}[H]
\centering
\includegraphics[width=0.85\textwidth]{fig_coefficient_plots.png}
\caption{Coefficient estimates from firm-level models A through E2. Bars indicate 95\% confidence intervals.}
\label{fig:firm_coefficients}
\end{figure}

\textbf{Model A} regresses the log of an automation proxy (constructed from capital expenditure on machinery and equipment) on log unit labor costs. The coefficient of 0.84 ($p < 0.001$, $R^2 = 0.43$) indicates that a 10\% increase in unit labor costs is associated with an 8.4\% increase in automation-related capital expenditure, consistent with a substitution elasticity meaningfully above unity. This is the central empirical result at the firm level: Colombian manufacturing firms respond to higher labor costs by investing in labor-saving capital, precisely as the theoretical framework predicts.

\textbf{Model C} confirms this pattern through an alternative specification: regressing log capital intensity on log labor cost per worker yields a coefficient of 1.27 ($p < 0.001$, $R^2 = 0.52$), implying that capital deepening responds more than proportionally to increases in per-worker labor costs. The elasticity exceeding unity is consistent with the substitution elasticities $\sigma > 1$ documented by \cite{cheng2021elasticity} using aggregate U.S. data.

\textbf{Model B} examines the investment rate (capital expenditure as a share of total assets) as a function of the labor share of value added. The positive coefficient of 0.045 ($p < 0.001$) indicates that firms with higher labor shares invest more intensively, consistent with the interpretation that high labor dependence creates stronger incentives for capital substitution.

\textbf{Models D and E2} reveal an apparent paradox. Model D, a logistic regression of EDIT-classified innovator status on unit labor costs, shows that higher ULC firms are \textit{less} likely to be classified as innovators. This finding reflects the structure of the EDIT survey, which classifies innovation broadly to include product innovation, organizational innovation, and marketing innovation---categories in which high-ULC firms (often small, labor-intensive establishments) are less active. However, Model E2, estimated on the subsample of firms that report \textit{any} positive investment in technology, shows that among firms that do invest, higher ULC is associated with substantially more intensive technology spending ($\beta = 0.81$, $p < 0.001$). The reconciliation is straightforward: high labor costs do not push all firms toward automation, but among firms that cross the investment threshold, cost pressure is a powerful determinant of investment intensity.

\begin{figure}[H]
\centering
\includegraphics[width=0.75\textwidth]{fig_scatter_lcpw_vs_invrate.png}
\caption{Scatter plot of labor cost per worker vs.\ investment rate across manufacturing firms. Each point represents one firm; the fitted line is from Model B.}
\label{fig:scatter_lcpw}
\end{figure}

The firm-size gradient in innovation rates is stark: only 9.9\% of micro firms (fewer than 10 employees) are classified as innovators in the EDIT, compared to 56.7\% of large firms (more than 200 employees). This gradient has direct implications for the distributional consequences of automation: large firms, which are disproportionately formal and capital-intensive, are far more likely to respond to cost increases by automating. Small firms, which lack the capital and technical capacity to automate, instead respond through informalization, wage compression, or exit---channels that reduce productivity and welfare without generating the efficiency gains that could offset displacement.

\begin{figure}[H]
\centering
\includegraphics[width=0.85\textwidth]{fig_sector_bubble_lcost_auto.png}
\caption{Bubble chart of manufacturing subsectors by labor cost per worker (horizontal axis), automation proxy (vertical axis), and employment (bubble size). The highest labor cost subsectors are coke/petroleum refining (60,932K COP/worker) and pharmaceuticals (54,758K COP/worker).}
\label{fig:sector_bubble}
\end{figure}

\subsection{Sectoral Vulnerability Analysis}
\label{subsec:sectoral_results}

We synthesize the individual-level, firm-level, and macroeconomic evidence into a composite Automation Vulnerability Index (AVI) that ranks Colombian economic sectors by their overall exposure to automation-driven displacement. The AVI integrates four dimensions: (i) the mean technical automation probability from the \cite{frey2017future} mapping; (ii) the labor cost incentive, measured as the ratio of total employer cost to base salary; (iii) the formality rate, which proxies the share of the sector's workforce subject to non-wage cost mandates; and (iv) a technology maturity indicator capturing the sector's current capital intensity and innovation rate.

\begin{figure}[H]
\centering
\includegraphics[width=0.85\textwidth]{fig_bar_vulnerability_ranking.png}
\caption{Automation Vulnerability Index (AVI) by economic sector. Higher values indicate greater composite vulnerability to automation-driven displacement.}
\label{fig:avi_ranking}
\end{figure}

Figure~\ref{fig:avi_ranking} displays the AVI rankings. Construction leads with a composite score of 80.5, driven by high automation probabilities for its predominantly routine physical tasks, high formality rates (construction firms are heavily regulated), and low current technology adoption implying substantial room for substitution. Administration, Education, and Health services follow at 78.0, a ranking that may appear counterintuitive given the low technical automation probabilities in education and health but reflects the high formality rates and labor cost incentives in public-sector employment, where parafiscal burdens are universally applied. Agriculture ranks third (72.0), combining high technical automation potential in routine crop and livestock tasks with an extremely low technology base, implying substantial vulnerability if precision agriculture and automated harvesting technologies continue to fall in cost.

Manufacturing (70.8) and Professional Services (70.5) occupy intermediate positions, while Financial Services (36.7), Utilities (40.8), and Real Estate (41.7) exhibit the lowest vulnerability scores. The low ranking of financial services---despite its high AI adoption rate---reflects the sector's combination of high worker skill levels, complex non-routine tasks, and the fact that current AI adoption in finance is largely augmentative (credit scoring, fraud detection) rather than substitutive.

\begin{figure}[H]
\centering
\includegraphics[width=0.85\textwidth]{fig_heatmap_sectoral_indicators.png}
\caption{Heatmap of sectoral indicators: TFP growth, labor share, formality rate, and automation risk. Darker shading indicates less favorable values.}
\label{fig:heatmap_sectors}
\end{figure}

The sectoral analysis reveals a troubling pattern when combined with productivity data. Total factor productivity (TFP) growth is negative in several of the most vulnerable sectors: Mining ($-2.14$\%), Financial and real estate services ($-1.06$\%), Construction ($-0.45$\%), and Manufacturing ($-0.43$\%). Negative TFP growth implies that these sectors are not merely stagnant but are experiencing technological regression---producing less output per unit of combined inputs over time. In this context, automation represents not an innovation-driven leap forward but a desperate attempt to arrest productivity decline, with the displacement of workers as a side effect rather than a byproduct of growth.

\begin{figure}[H]
\centering
\includegraphics[width=0.85\textwidth]{fig_scatter_productivity_vs_laborshare.png}
\caption{Scatter plot of labor productivity growth vs.\ labor share of value added by sector. Sectors in the upper-left quadrant combine high labor dependence with low productivity growth---the zone of maximum automation vulnerability.}
\label{fig:prod_vs_laborshare}
\end{figure}

\subsection{Generative AI Impact}
\label{subsec:genai}

The results presented above are based primarily on the \cite{frey2017future} framework, which maps automation probabilities to occupations based on the technical feasibility of automating their constituent tasks using pre-2017 technologies---predominantly industrial robotics, rule-based software, and early machine learning. The emergence of generative AI (GenAI) since 2022, and particularly the deployment of large language models capable of performing complex cognitive tasks, has fundamentally altered the automation risk landscape in ways that the Frey--Osborne framework does not capture.

The World Economic Forum's \textit{Future of Jobs Report 2025} \cite{wef2025} identifies generative AI as the single most transformative technology for labor markets over the 2025--2030 horizon, with 86\% of surveyed employers expecting it to transform their business processes. The OECD/GPAI report on Latin America \cite{oecd2025voces} documents that GenAI disproportionately affects cognitive, non-routine tasks---precisely the tasks that the traditional framework classifies as low-risk. This creates what has been termed a ``reverse skill bias'': whereas previous waves of automation displaced routine manual and routine cognitive workers (the middle of the skill distribution), GenAI threatens non-routine cognitive workers (the upper-middle and upper segments), including accountants, lawyers, translators, content creators, financial analysts, and software developers.

For Colombia, the GenAI channel is particularly consequential for two reasons. First, the BPO sector---which employs over 600,000 workers and has been a major source of formal job creation---is heavily exposed. BPO tasks (customer service, data entry, translation, content moderation) are precisely the cognitive routine tasks that large language models can perform at a fraction of the human labor cost. The 80\% AI adoption rate already reported in the sector suggests that this displacement is not prospective but ongoing. Second, the financial sector's aggressive AI adoption, while currently augmentative, is transitioning toward substitutive applications: Bancolombia's investment in AI-assisted coding (4,000 tech workers, GitHub Copilot deployment) is designed to increase per-developer productivity, which arithmetically implies fewer developers needed for a given volume of output.

The ILIA 2024 \cite{ilia2024} ranking---Colombia fifth in Latin America with a score of 52.64, behind Chile (73.07), Brazil (69.30), and Uruguay (64.98)---suggests that Colombia's capacity to manage the GenAI transition lags its exposure to it. The gap relative to Chile is driven by differences in governance frameworks, public data infrastructure, and the density of university-industry AI research partnerships. Colombia's CONPES 4144 \cite{conpes4144} represents an attempt to close this gap, but as we discuss in Section~\ref{sec:policy}, its resource allocation and institutional design fall short of what the scale of the challenge requires.


% ============================================================================
\section{Policy Landscape and Regulatory Paradox}
\label{sec:policy}
% ============================================================================

Colombia's policy response to automation risk is characterized by a fundamental tension: the same regulatory framework that generates high formal labor costs---thereby incentivizing automation---simultaneously funds the institutions (SENA, the health and pension systems) that are supposed to protect displaced workers. This section examines the three principal policy domains---labor regulation, AI governance, and workforce development---and argues that their current configuration is not merely suboptimal but internally contradictory.

\subsection{The Archived Labor Reform}
\label{subsec:labor_reform}

The labor reform bill debated and ultimately archived by the Colombian Congress in March 2025 proposed a suite of measures that would have substantially increased employer costs. The principal provisions included: extending the night-shift premium to begin at 7:00 PM (from the current 9:00 PM), raising the Sunday and holiday premium to 100\% of base salary, restricting the use of fixed-term and temporary contracts, expanding severance protections, and mandating additional rest periods. Business associations, led by Fenalco, estimated that the combined effect of these provisions would increase total labor costs by 18--34\%, depending on the sector and the share of non-standard hours in the production schedule.

The automation implications of such a cost increase are direct and quantifiable. In employer surveys conducted during the legislative debate, 17\% of respondents indicated that they would respond to the reform by substituting technology for labor. While this figure may appear modest, it represents the \textit{declared} intention of firms surveyed in a politically charged environment; the actual substitution response, unfolding gradually over years as capital investment decisions are made, would likely exceed the survey estimate. The international evidence reviewed in Section~\ref{sec:international}---particularly Korea's experience with a 16.4\% minimum wage increase---suggests that the 18--34\% cost increase proposed by the reform falls squarely within the range that triggers measurable automation responses.

Although the reform was archived, its provisions remain part of the policy repertoire and may be reintroduced. More importantly, the reform debate revealed a structural feature of Colombian labor policy: proposed labor protections are designed as if automation were not an available margin of adjustment. The implicit model is one in which firms can only respond to cost increases by accepting lower profits or raising prices---a model that ignores the substitution elasticities documented in Section~\ref{subsec:panel_results} and Section~\ref{subsec:firm_results}.

\subsection{CONPES 4144 and National AI Policy}
\label{subsec:conpes}

CONPES 4144 \cite{conpes4144}, Colombia's national AI policy approved in 2025, allocates COP 479 billion (approximately USD 110 million) through 2030 across 106 discrete actions organized around five strategic pillars: AI governance, talent development, research and innovation, data infrastructure, and sectoral adoption. The policy represents a significant improvement over previous ad hoc initiatives and correctly identifies many of the relevant challenges. However, several limitations constrain its potential impact.

First, the resource envelope is modest relative to the scale of the challenge. COP 479 billion over five years amounts to approximately COP 96 billion per year---less than 0.006\% of GDP. By comparison, Korea's AI investment plan allocates over USD 2 billion annually, and even Chile's more modest national AI strategy is backed by substantially higher per-capita funding. Second, the 106 actions are distributed across multiple ministries and agencies with no single coordinating authority empowered to enforce implementation timelines, creating a risk of fragmentation and bureaucratic dilution.

The workforce development component of CONPES 4144 merits particular scrutiny. SENATEC, the digital skills training initiative operated through SENA, has trained approximately 303,000 people in IT-related skills---a figure that represents only 1.3\% of the total workforce. Moreover, the quality and labor market relevance of these training programs remain unverified: SENA's curriculum development processes are slow, and the institution's IT training offerings have historically lagged several years behind industry requirements. The gap between the 303,000 workers trained in basic IT skills and the millions facing automation risk documented in Section~\ref{subsec:risk_colombia} is indicative of the mismatch between the scale of the problem and the scale of the policy response.

\subsection{Institutional Disconnect}
\label{subsec:disconnect}

The deepest structural problem in Colombia's policy architecture is the paradoxical relationship between parafiscal contributions and workforce adaptation. SENA---the institution responsible for vocational training and reskilling programs that could, in principle, equip workers for the automated economy---is funded primarily by a 2\% payroll tax levied on formal employers. This creates a perverse incentive loop: the parafiscal contribution that funds SENA simultaneously increases the cost of formal labor, thereby incentivizing automation, which displaces the formal workers whose payroll taxes fund SENA. As automation reduces formal employment, SENA's funding base erodes, diminishing its capacity to retrain displaced workers precisely when the need for retraining is greatest.

This ``SENA paradox'' is compounded by the broader parafiscal structure. The 2012 tax reform (Law 1607) eliminated employer contributions to SENA, ICBF, and the health system for workers earning less than 10 minimum wages, replacing these revenues with a general tax (CREE, later integrated into the income tax). However, employer pension contributions (12\%), health contributions for higher earners (8.5\%), and other mandatory charges remain substantial, and the net non-wage cost burden continues to exceed 46\% of base salary for the average formal worker.

The jornada laboral reduction mandated by Law 2101 of 2021---which progressively reduces the standard workweek from 48 to 42 hours, with full implementation by July 2026---adds another cost dimension. Unless accompanied by proportional productivity increases, the reduction in hours at unchanged monthly pay effectively raises the hourly labor cost by approximately 14\%. While the law is designed to improve work-life balance, its effect on the automation calculus is unambiguous: every increase in the effective hourly cost of labor shifts the cost-benefit frontier in favor of technological substitution.

Colombia's institutional architecture thus lacks the three elements that have enabled other countries to manage automation transitions successfully: (i) a dual vocational training system, as in Germany, that continuously adapts worker skills to evolving technological requirements; (ii) an effective social dialogue mechanism that gives workers a voice in firm-level automation decisions; and (iii) sufficient R\&D investment to ensure that automation proceeds through productivity-enhancing channels rather than pure cost-driven substitution. Without these elements, automation in Colombia is likely to follow a ``low road'' trajectory: firms adopt technologies that are merely cheaper than formal labor rather than genuinely more productive, generating displacement without the compensating welfare gains that would make automation socially beneficial.


% ============================================================================
\section{Scenario Projections}
\label{sec:scenarios}
% ============================================================================

To quantify the potential employment consequences of automation under alternative policy and technological trajectories, we construct a scenario simulation model that projects formal employment displacement over a ten-year horizon (2025--2035). The model combines the sectoral automation risk estimates from Section~\ref{subsec:risk_colombia}, the firm-level substitution elasticities from Section~\ref{subsec:firm_results}, and the international adoption trajectories from Section~\ref{subsec:panel_results} to generate displacement estimates under five scenarios that span the plausible range of outcomes.

\subsection{Scenario Definitions}
\label{subsec:scenario_definitions}

The five scenarios are defined as follows:

\begin{enumerate}[label=(\roman*)]
\item \textbf{Status Quo:} Current regulatory framework, current automation adoption trajectory, no major policy changes. Labor costs grow in line with the historical trend of minimum wage increases (approximately 3--4\% real per year). AI adoption continues at the pace established in 2023--2025.

\item \textbf{Labor Reform:} The archived labor reform is reintroduced and enacted, increasing total labor costs by 18--34\% (we use the midpoint of 26\%). All other parameters follow the Status Quo trajectory.

\item \textbf{Parafiscal Reform:} Parafiscal contributions are reduced or eliminated for formal employment, lowering the non-wage cost wedge from 52\% to approximately 30\% of base salary. This scenario represents the ``German path'' of reducing labor costs to slow the automation incentive.

\item \textbf{AI Acceleration:} Generative AI capabilities advance faster than current baseline projections, reducing the cost of cognitive automation by 50\% relative to the Status Quo. This scenario captures the ``GPT-5 and beyond'' trajectory in which large language models achieve near-human performance on a wider range of cognitive tasks.

\item \textbf{Combined Worst Case:} The Labor Reform is enacted \textit{simultaneously} with the AI Acceleration scenario, representing the convergence of maximum cost pressure with maximum technological capability.
\end{enumerate}

\subsection{Simulation Results}
\label{subsec:simulation_results}

Figure~\ref{fig:trajectories} displays the projected formal employment trajectories under each scenario.

\begin{figure}[H]
\centering
\includegraphics[width=0.90\textwidth]{fig_sim_formal_employment_trajectories.png}
\caption{Projected formal employment trajectories under five scenarios, 2025--2035. Shaded bands represent 95\% confidence intervals from Monte Carlo simulation (10,000 draws).}
\label{fig:trajectories}
\end{figure}

\begin{table}[H]
\centering
\caption{Scenario Projection Summary: Formal Employment Displacement by 2035}
\label{tab:scenarios}
\begin{threeparttable}
\begin{tabular}{lccc}
\toprule
\textbf{Scenario} & \textbf{Formal Empl.\ Change} & \textbf{Displaced Workers} & \textbf{95\% CI} \\
\midrule
Status Quo & $-18.9\%$ & 3.34M & [3.21M, 3.54M] \\
Labor Reform & $-39.0\%$ & 4.32M & [4.19M, 4.50M] \\
Parafiscal Reform & $-14.6\%$ & 2.80M & [2.62M, 3.11M] \\
AI Acceleration & $-25.5\%$ & 3.85M & [3.76M, 3.98M] \\
Combined Worst & $-42.7\%$ & 4.31M & [4.22M, 4.45M] \\
\bottomrule
\end{tabular}
\begin{tablenotes}
\small
\item \textit{Notes:} Displaced workers represent the cumulative number of formal positions eliminated over the 2025--2035 horizon. Confidence intervals derived from Monte Carlo simulation with 10,000 draws, varying automation probabilities, substitution elasticities, and adoption rates within their estimated distributions.
\end{tablenotes}
\end{threeparttable}
\end{table}

The results warrant several observations. Under the \textbf{Status Quo}, automation displaces an estimated 3.34 million formal workers over the decade, representing a 18.9\% decline in formal employment. This is not a benign baseline: it reflects the continued operation of existing cost incentives, the ongoing diffusion of AI in services, and the progressive reduction of robot costs. The displacement is concentrated in manufacturing, transport, and formal retail, with professional services and BPO experiencing accelerating losses as generative AI matures.

The \textbf{Labor Reform} scenario amplifies displacement dramatically. The 18--34\% increase in labor costs pushes an additional $\sim$1 million formal positions past the automation threshold, bringing total displacement to 4.32 million ($-39.0\%$). Figure~\ref{fig:waterfall} decomposes this increment: the single largest contributor is the extension of the night-shift premium, which disproportionately affects manufacturing and logistics firms operating multiple shifts.

\begin{figure}[H]
\centering
\includegraphics[width=0.85\textwidth]{fig_sim_waterfall_labor_reform.png}
\caption{Waterfall decomposition of incremental displacement under the Labor Reform scenario relative to Status Quo. Each bar represents the contribution of a specific reform provision.}
\label{fig:waterfall}
\end{figure}

The \textbf{Parafiscal Reform} scenario demonstrates the substantial potential of cost-side interventions. By reducing the non-wage cost wedge, this scenario limits displacement to 2.80 million workers ($-14.6\%$)---saving approximately 1.5 million formal jobs relative to the Labor Reform scenario and 540,000 relative to the Status Quo. This finding constitutes the most actionable policy result of the study: parafiscal reform is the single most powerful lever available to Colombian policymakers for slowing automation-driven displacement, because it directly reduces the cost differential that makes formal workers more expensive than their technological substitutes.

Under the \textbf{AI Acceleration} scenario, displacement reaches 3.85 million ($-25.5\%$), reflecting the extension of automation risk to cognitive, non-routine occupations in professional services, BPO, and financial services. The \textbf{Combined Worst Case} produces displacement of 4.31 million ($-42.7\%$), a figure nearly identical to the Labor Reform scenario alone, suggesting that the Labor Reform's cost increase is so large that it dominates even aggressive AI acceleration in determining the total displacement envelope.

\begin{figure}[H]
\centering
\includegraphics[width=0.85\textwidth]{fig_sim_heatmap_displacement.png}
\caption{Heatmap of projected displacement by sector and scenario. Cell values represent percentage decline in formal employment.}
\label{fig:heatmap_displacement}
\end{figure}

The sectoral decomposition (Figure~\ref{fig:heatmap_displacement}) reveals extreme heterogeneity. In the Combined Worst Case, BPO and Professional Services lose 57--65\% of formal positions, reflecting the dual pressure of high labor costs and high cognitive automation exposure. Manufacturing loses 35--45\%, driven primarily by the physical automation channel. Agriculture and construction experience more moderate losses (20--30\%) owing to lower formality rates, which insulate a large share of the workforce from cost-driven substitution.

\subsection{Sensitivity Analysis}
\label{subsec:sensitivity}

We assess the sensitivity of the displacement estimates to perturbations in the key input parameters through a tornado analysis (Figure~\ref{fig:tornado}).

\begin{figure}[H]
\centering
\includegraphics[width=0.85\textwidth]{fig_sim_tornado_sensitivity.png}
\caption{Tornado chart of sensitivity analysis for the Status Quo scenario. Each bar shows the change in total displacement when the indicated parameter is varied by $\pm 1$ standard deviation.}
\label{fig:tornado}
\end{figure}

The most sensitive parameter is the mean automation probability: a $\pm 10$ percentage point shift in the automation risk distribution generates a range of 1.65 million workers in the displacement estimate. This sensitivity underscores the importance of the mapping methodology used to assign automation probabilities to Colombian occupations. To the extent that the Frey--Osborne framework---developed for U.S. occupations and mapped to Colombian categories at the one-digit level---overstates or understates the true automation potential of Colombian jobs, the displacement estimates will shift proportionally.

The second most sensitive parameter is the substitution elasticity, followed by the discount rate applied to future automation cost reductions. The policy cost increase (relevant only in the Labor Reform and Combined Worst Case scenarios) ranks fourth, suggesting that while policy matters, the underlying technological trajectory and labor market structure are more fundamental determinants of the displacement envelope.

\begin{figure}[H]
\centering
\includegraphics[width=0.85\textwidth]{fig_sim_fan_chart_baseline.png}
\caption{Fan chart of the Status Quo displacement trajectory showing the 50\%, 75\%, and 95\% prediction intervals from Monte Carlo simulation.}
\label{fig:fan_chart}
\end{figure}

Figure~\ref{fig:fan_chart} presents the fan chart for the Status Quo scenario, illustrating the widening uncertainty band over the projection horizon. By 2030, the 95\% prediction interval spans approximately 1.5 million workers; by 2035, it exceeds 2.5 million. This widening reflects the compounding of uncertainty in technological progress, policy responses, and macroeconomic conditions over a decade-long horizon, and counsels against treating the point estimates as precise forecasts. The value of the scenario exercise lies not in the specific numbers but in the comparative magnitudes: the robust finding is that parafiscal reform substantially reduces displacement relative to the Status Quo, while labor cost increases substantially amplify it.


% ============================================================================
\section{Conclusion}
\label{sec:conclusion}
% ============================================================================

This study has investigated the relationship between labor costs and automation in Colombia through the lens of a multi-level empirical framework spanning international panels, individual-level microdata, firm-level surveys, and sectoral aggregates. Five principal findings emerge.

\textbf{First}, international panel regressions confirm that labor productivity and labor cost structures are robust determinants of robot adoption across 25 countries, with strong persistence in automation trajectories. Colombia installs roughly one-tenth the number of robots predicted by its macroeconomic fundamentals, a ``missing robots'' gap attributable to small firm size, weak supply-chain integration, limited capital access, and the availability of cheap informal labor.

\textbf{Second}, approximately 53.8\% of Colombian workers occupy occupations classified as high automation risk ($\Pr \geq 0.50$), with formal workers facing 3.52 times the odds of high-risk employment relative to informal workers. This formality paradox---whereby the regulatory framework designed to protect workers simultaneously makes them more expensive to employ and therefore more attractive to replace with technology---is the defining structural feature of Colombia's automation challenge.

\textbf{Third}, firm-level evidence from the matched EAM--EDIT dataset confirms that Colombian manufacturing firms respond to higher labor costs by increasing capital intensity, with an estimated substitution elasticity of 0.84 between unit labor costs and automation-related investment. Among firms that invest in technology, the intensity of investment rises monotonically with labor cost pressure ($\beta = 0.81$, $p < 0.001$), confirming the cost-driven adoption mechanism at the establishment level.

\textbf{Fourth}, the Automation Vulnerability Index identifies Construction (80.5), Administrative/Education/Health services (78.0), Agriculture (72.0), and Manufacturing (70.8) as the most vulnerable sectors, while Financial Services (36.7) and Utilities (40.8) exhibit the lowest composite vulnerability---though financial services faces accelerating cognitive automation risk through the generative AI channel.

\textbf{Fifth}, scenario projections indicate that the archived labor reform, if enacted, would displace an additional 1.0 million formal workers relative to the Status Quo baseline (4.32M vs.\ 3.34M), while a parafiscal reform reducing non-wage costs would save approximately 540,000 formal positions (2.80M vs.\ 3.34M). The differential between the reform and parafiscal scenarios---approximately 1.5 million formal jobs---represents the stakes of the policy choice facing Colombian legislators.

These findings carry three direct policy implications. \textbf{The most urgent is parafiscal reform.} Reducing the non-wage cost wedge that currently inflates the total employer cost of formal labor by 46--58\% above base salary would simultaneously strengthen formal employment, reduce the automation incentive, and improve Colombia's competitiveness relative to regional peers. The 2012 reform (Law 1607) demonstrated that parafiscal reduction is politically feasible; extending it to cover remaining contributions---while compensating SENA and other beneficiary institutions through general taxation---should be a legislative priority.

\textbf{Second, Colombia must invest substantially in reskilling infrastructure.} The current SENATEC program, at 303,000 trainees representing 1.3\% of the workforce, is orders of magnitude below the scale required. The German model of dual vocational training, which continuously adapts curricula to evolving industry requirements through employer-union co-governance of training content, provides a template. However, transplanting this model requires institutional capacities---effective employer associations, representative unions, and sustained public investment in training infrastructure---that Colombia currently lacks and must build.

\textbf{Third, policymakers must avoid large, discrete labor cost shocks.} The international evidence---particularly from Korea's 2018 minimum wage increase---demonstrates that abrupt cost increases trigger non-linear automation responses, producing waves of simultaneous displacement that overwhelm the absorptive capacity of other sectors and of social protection systems. If labor protections are to be strengthened, they should be phased gradually and accompanied by offsetting reductions in non-wage costs, so that the total employer cost trajectory remains smooth and predictable.

The study is subject to several limitations that warrant explicit acknowledgment. The individual-level analysis relies on cross-sectional GEIH data, which precludes the identification of causal effects or the tracking of individual workers over time. The EAM--EDIT merge involves a temporal mismatch, as the two surveys do not share identical reference periods, introducing measurement error in the firm-level estimates. The mapping of Frey--Osborne automation probabilities to Colombian occupations is performed at the one-digit level of the occupational classification, which masks substantial within-category heterogeneity. And the scenario projections, while grounded in estimated parameters, are inherently speculative over a ten-year horizon and should be interpreted as illustrative of magnitudes rather than precise forecasts.

Despite these limitations, the central message is robust across specifications and analytical levels: Colombia's high non-wage labor costs create a structural incentive for automation that is amplified by regulatory proposals to increase those costs further. The window for proactive policy response is narrowing. The emergence of generative AI has compressed the timeline for cognitive automation by approximately a decade relative to pre-2022 projections \cite{wef2025}, meaning that the displacement pressures historically associated with industrial robots---which took decades to materialize in manufacturing---may unfold over years in services. Colombia's challenge is not whether automation will transform its labor market, but whether the transformation will be managed through deliberate institutional adaptation or suffered as a disorderly shock to an unprepared workforce. The evidence presented in this paper suggests that, absent significant policy reform, the latter outcome is the more probable.